\documentclass{article}
\usepackage{amsmath,amssymb}
\usepackage{xurl}
\usepackage[hidelinks]{hyperref}
% Shared commands for notes in docs/paper-gaps/.

\DeclareUnicodeCharacter{2080}{\ifmmode{}_{0}\else\textsubscript{0}\fi}
\DeclareUnicodeCharacter{2081}{\ifmmode{}_{1}\else\textsubscript{1}\fi}
\DeclareUnicodeCharacter{2082}{\ifmmode{}_{2}\else\textsubscript{2}\fi}
\DeclareUnicodeCharacter{2099}{\ifmmode{}_{n}\else\textsubscript{n}\fi}

\newcommand{\Lean}{\textsc{Lean}}
\ExplSyntaxOn
\NewDocumentCommand{\leanid}{m}{
  \ifmmode
    \textsf{\detokenize{#1}}
  \else
    \group_begin:
    \urlstyle{sf}
    \tl_set:Nn \l_tmpa_tl {#1}
    \tl_replace_all:Nnn \l_tmpa_tl {\_} {_}
    \exp_args:NV \path \l_tmpa_tl
    \group_end:
  \fi
}
\ExplSyntaxOff
\providecommand{\tr}{\operatorname{tr}}
\newcommand{\tnleanIssueBase}{https://github.com/LionSR/TNLean/issues/}
\newcommand{\tnleanPrBase}{https://github.com/LionSR/TNLean/pull/}
\newcommand{\tnleanDocBase}{https://sirui-lu.com/TNLean/docs/}
\newcommand{\ghissue}[1]{\href{\tnleanIssueBase#1}{GitHub issue \##1}}
\newcommand{\ghpr}[1]{\href{\tnleanPrBase#1}{PR \##1}}
\newcommand{\doclink}[1]{\href{\tnleanDocBase#1.html}{\path{#1}}}

% Dirac notation and the identity operator, as in blueprint/src/macros/common.tex.
% \providecommand keeps note-local definitions authoritative.
\usepackage{dsfont}
\providecommand{\ket}[1]{|#1\rangle}
\providecommand{\bra}[1]{\langle #1|}
\providecommand{\braket}[2]{\langle #1 | #2 \rangle}
\providecommand{\ketbra}[2]{|#1\rangle\!\langle #2|}
\providecommand{\Id}{\mathds{1}}
\providecommand{\one}{\mathds{1}}
\providecommand{\C}{\mathbb{C}}
\providecommand{\MN}[1]{M_{#1}(\C)}
\providecommand{\MD}{M_D(\C)}

% External referencing for paper sources.  The build helper writes sanitized
% auxiliary files under build/paper-aux/ and excludes duplicated source labels.
\usepackage{xr}

% Import a generated paper auxiliary file with a caller-supplied label prefix.
% The candidate paths support builds from docs/paper-gaps/, the repository root,
% and build/paper-aux/ itself.
\newcommand{\importpaperaux}[2]{%
  \IfFileExists{../../build/paper-aux/#2.aux}{%
    \externaldocument[#1]{../../build/paper-aux/#2}%
  }{%
    \IfFileExists{build/paper-aux/#2.aux}{%
      \externaldocument[#1]{build/paper-aux/#2}%
    }{%
      \IfFileExists{#2.aux}{%
        \externaldocument[#1]{#2}%
      }{}%
    }%
  }%
}

% General external reference macros
\newcommand{\paperref}[2]{\ref{#1-#2}}
\newcommand{\papereqref}[2]{(\ref{#1-#2})}

% CPSV16 source (1606.00608 / MPDO-22-12-17-2.tex) external document and shortcuts
\importpaperaux{cpsv16-}{1606.00608/MPDO-22-12-17-2-tnlean}
\newcommand{\cpsvref}[1]{\paperref{cpsv16}{#1}}
\newcommand{\cpsveqref}[1]{\papereqref{cpsv16}{#1}}

% Verdict marker: \gapnote{<kind>}{<status>}. Kind is the policy
% classification (clarification, local-correction, scope-restriction,
% unfaithful, false-source, open-gap); status is open, wip (elimination
% underway), resolved, or historical. Machine-read by the site index;
% prints a verdict line.
\newcommand{\gapnote}[2]{\par\noindent\textbf{Verdict:} #1 (#2).\par}

\title{The CZX Tensor: Orientation of the Bond Labels}
\author{}
\date{2026-09-26}
\begin{document}
\maketitle
\gapnote{local-correction}{open}

\section{What the source prints}
Section~``MPS and PEPS'' of \cite{Cirac2021Matrix} contracts a PEPS by
placing on every edge the maximally entangled pair
$\sum_{n}\lvert n)\otimes\lvert n)$, so the two ends of a bond carry the same
label.\footnote{Source traceability:
    \path{Papers/2011.12127/TN-Review-main.tex}, lines 271--276.}
In Appendix~A, under ``Two dimensions: PEPS'', the virtual indices of a site
tensor are ordered top, right, down and left.\footnote{Source traceability:
    \path{Papers/2011.12127/TN-Review-main.tex}, line 2415.}
The CZX state of \cite{ChenLiuWen2011CZX} is the product over plaquettes of
the four-qubit states $\ket{0000}+\ket{1111}$, and the review writes it,
after blocking the four qubits of a site, as the PEPS with tensor
\begin{align}
    \sum_{i,j,k,l=0}^{1}\ket{ijkl}\bigl((i,j),(j,k),(k,l),(l,i)\bigr|,
    \label{eq:czx_printed}
\end{align}
where $i,j,k,l$ are the qubits at the top-left, top-right, bottom-right and
bottom-left corners of the site.\footnote{Source traceability:
    \path{Papers/2011.12127/TN-Review-main.tex}, lines 2502--2514;
    \path{References/1106.4752/source/dDSPTmodel.tex}, lines 317--321.}

\section{The literal contraction}
Write $(i_v,j_v,k_v,l_v)$ for the qubits of the site $v$, and $v+e_1$,
$v+e_2$ for its right and upper neighbours. Contracting
\eqref{eq:czx_printed} with equal labels at the two ends of every bond, the
right leg $(j_v,k_v)$ of $v$ meets the left leg $(l_{v+e_1},i_{v+e_1})$ of
$v+e_1$, and the top leg $(i_v,j_v)$ of $v$ meets the down leg
$(k_{v+e_2},l_{v+e_2})$ of $v+e_2$. The nonzero coefficients are those with
\begin{align}
    j_v=l_{v+e_1}=l_{v+e_2},\qquad k_v=i_{v+e_1},\qquad i_v=k_{v+e_2}.
    \notag
\end{align}
These conditions join the top-right qubit of $v$ to bottom-left qubits of
two different sites, and they chain the qubits along the diagonals
$v\mapsto v+e_1-e_2$ and $v\mapsto v+e_1+e_2$. The state is a product of
GHZ states on diagonal loops that wind around the torus, not the product of
plaquette states. On a torus of $3\times 3$ sites, for instance, a plaquette
configuration with a single plaquette flipped has coefficient $0$ in the
literal contraction. The failure is on nondegenerate data.

\section{The correction}
Two neighbouring sites list the two qubits of their common side in opposite
orders, clockwise around each site. The plaquette state is obtained when a
bond identifies the pair $(a,b)$ at one end with $(b,a)$ at the other, so
that each bond joins the two qubits of one plaquette to the two qubits of the
same plaquette across the common side, as in the figure of the blocked state
in \cite{Cirac2021Matrix}. Equivalently, the tensor \eqref{eq:czx_printed} is used with its
down and left labels reversed. With this reading the contraction gives
\begin{align}
    c(\sigma)=\prod_{v}\delta\bigl[j_v=i_{v+e_1}=k_{v+e_2}
    =l_{v+e_1+e_2}\bigr],
    \notag
\end{align}
the product over plaquettes of the unnormalized states
$\ket{0000}+\ket{1111}$, on every torus of width and height at least
three.\footnote{Results: \leanid{TNLean.PEPS.czxPEPS},
    \leanid{TNLean.PEPS.stateCoeff_czxPEPS}.}

\section{Elimination plan}
None needed beyond this record: the printed tensor is kept, and the reversal
of the down and left labels is part of the definition of the PEPS. A reading
of the source with the reversed identification on every bond would state the
same theorem.

\bibliographystyle{alpha}
\bibliography{references}
\end{document}
